\chapter{Downloading public data through APIs}\label{ch:apis}

The funder wants a line in the report about how your program's counties compare to the state on child poverty. The number exists at \texttt{census.gov}, and you have twenty minutes before the draft is due. An application programming interface, or API\index{API (application programming interface)}, is the door that makes twenty minutes enough: a web address you can query from Stata and get data back, no browser and no copy-paste.

This chapter teaches the grammar of that door once, then the vocabulary of the agencies worth knowing: Census, education, economic, and health data, each reachable from Stata with a package, a native command, or a few lines of Python. The worked example in Section~\ref{sec:qcew} downloads real county wage data and builds the figure on the facing page, start to finish, with no key. The point throughout is that a scripted download is a replicable download: the number in the report rebuilds next year by rerunning the same lines, which a screenshot never can.

\section{Practitioners hold the questions; public APIs hold the denominators}\label{ch03-advantage}
Before the mechanics, it is worth saying plainly why this chapter's skill belongs near the front of an embedded evaluator's toolkit. The program team you sit with owns the numerators: enrollments, completions, placements, the counts their information system produces every week. What that system cannot produce is the denominator, the county population, the district enrollment, the regional wage, that turns a count into a rate a board can judge. Performance measurement runs on exactly this conversion; a logic model's outcomes column names community-level change, and community-level change is only visible against a community-level baseline. The embedded evaluator's comparative advantage is sitting close enough to the practitioners to know which denominator matters, while owning the scripted route to fetch it.

The research literature has been moving toward these same sources for over a decade, so an evaluator who leans on them is standing on firm ground. Empirical economics shifted from small survey samples toward administrative records with near-universal coverage, a shift \textcite{einav2014bigdata} describe as the defining data development of the era, and one that \textcite{card2010administrative} argued for on quality grounds: records generated by the administration of taxes and programs cover more people, more accurately, than any survey an evaluator could field. The push has a mirror image on the survey side, where \textcite{meyer2015surveys} document rising nonresponse, imputation, and measurement error in the household surveys evaluators once treated as the default benchmark. The so-what for a nonprofit analyst is concrete: the QCEW wage file and the CCD enrollment file in this chapter are administrative censuses, not samples, and citing them puts your context numbers on the same footing as the published literature's. The turn is not confined to economics: a full RSF volume surveys administrative data across sociology, education, and criminology, and \textcite{penner2019administrative} frame its value for policy in terms an embedded evaluator will recognize, that records generated by service delivery answer questions surveys cannot afford to and track the same people over time without re-contacting them.

Administrative data earn one standing caution, which this chapter's margin notes keep repeating in specific forms. These records are by-products of running a program, not instruments designed for research, so their definitions follow program rules rather than research questions \parencite{connelly2016administrative}: the QCEW counts insured jobs rather than employed residents, and the CCD counts fall-snapshot enrollment rather than students served. Knowing the generating process behind each file is the evaluator's half of the bargain, and it is the same habit of interrogating measures that performance-measurement practice already teaches.

The reason these files sit behind open endpoints at all is policy, not accident. The 2009 open-government memorandum directed federal agencies toward transparency, participation, and collaboration \parencite{obama2009transparency}, and the \texttt{data.gov} catalog and the agency APIs of this chapter are that directive's working infrastructure, later reinforced in statute by the Foundations for Evidence-Based Policymaking Act, which made open, machine-readable data and agency evidence-building a legal default rather than an aspiration \parencite{congress2019evidence}. That institutional footing matters to an evaluator in two practical ways: the endpoints are citable public assets rather than favors from a webmaster, and they are likelier to still answer next year, which is what a replicable download quietly depends on.

Timing closes the case. Utilization-focused evaluation's central claim is that findings get used when they arrive inside the decision window of an intended user \parencite{patton2008utilization}, and the twenty-minute funder request that opened this chapter is that claim in miniature. A context number that takes three weeks of data requests misses the board meeting; the same number pulled by a four-line script makes the draft. The same speed serves designs, not just deadlines: when a comparison-group study needs to show that program counties resembled the rest of the state at baseline, the kind of baseline-equivalence evidence the What Works Clearinghouse standards require of non-randomized designs before a study can meet standards \parencite{wwc2022handbook}, the baseline series comes from these same public endpoints.

\section{The shape of an API request}
Learn the anatomy of one API request and you can read every portal in this chapter, because underneath they are the same object: a base URL, a path that names the dataset, and a query string of parameters that names the slice you want. Some return comma-separated text you can read directly; others return JSON, a nested text format that needs one parsing step. Many require a free key, a short string that identifies you to the server. The key belongs in your \texttt{profile.do} as a global macro, never pasted into shared code, because a key in a public repository is a key someone else will spend. Figure~\ref{fig:ch03-viz-anatomy} maps those three parts onto the actual QCEW URL the worked example downloads.

\begin{figure*}[htbp]
\centering
\begin{tikzpicture}[node distance=4mm]
% the three parts of any request, left to right, then the reply
\node[stage] (base) {base URL};
\node[stage, right=of base] (path) {path};
\node[stage, right=of path] (query) {query\\string};
\node[io, right=8mm of query] (reply) {CSV or\\JSON reply};
\draw[flow] (query) -- (reply);
% each part tied to a piece of the real QCEW URL
\node[substep, below=5mm of base, text width=2.1cm] (baseA)
  {\texttt{data.bls.gov}\\the server};
\node[substep, below=5mm of path, text width=2.6cm] (pathA)
  {\texttt{/cew/data/api}\\names the dataset};
\node[substep, below=5mm of query, text width=2.6cm] (queryA)
  {\texttt{2023/1/area/}\\ \texttt{48453.csv}\\the slice you want};
\draw[flow, dashed] (base) -- (baseA);
\draw[flow, dashed] (path) -- (pathA);
\draw[flow, dashed] (query) -- (queryA);
\node[annot, below=4mm of pathA.south, xshift=6mm] (key)
  {some endpoints add a free \emph{key} identifying you to the server};
\end{tikzpicture}
\caption[Anatomy of an API request]{Every portal in this chapter is the same object: a base URL, a path that names the dataset, and a query string that names the slice, returning CSV or JSON. Read the labels bottom-up against the QCEW call \texttt{data.bls.gov/cew/data/api/2023/1/area/48453.csv}: swap the FIPS code in the query string and you have a different county.}
\label{fig:ch03-viz-anatomy}
\end{figure*}

Learning the pattern rather than memorizing one agency is what makes this skill extensible: the next portal you meet has a different path and different parameters, but the same anatomy. We also request politely, spacing calls and taking only the rows we need, because an embedded evaluator depends on these public services continuing to answer.\marginnote{Rate-limiting etiquette: a \texttt{sleep 500} between calls costs you half a second and costs the server nothing. Hammer an endpoint with a tight loop and you may earn an HTTP 429, or a quiet IP ban that outlasts your deadline.}

\subsection{Rate-limit etiquette is commons stewardship, not just self-protection}\label{ch03-commons}
The polite pause between calls is easy to read as a trick for dodging an error code, but it is better understood as the evaluator's small share of maintaining a shared resource. A public API is a commons: no single analyst pays for the server, every analyst depends on it, and the failure mode is congestion when many users hit it hard at once. The evaluator who caps a loop with \texttt{sleep} and requests only the county-years needed is doing what commons stewardship always asks, taking a sustainable slice so the resource is still there for the next person, who may be a colleague at another nonprofit running the same query next week. Framing it this way turns an annoyance into a professional obligation, and it is the same obligation the field already recognizes: the AEA Guiding Principles ask evaluators to act with integrity and regard for the public welfare beyond the immediate engagement \parencite{aea2018principles}, and courtesy toward the public infrastructure other evaluators rely on is that principle at the level of a do-file.

The stewardship habit has three concrete moves, none of them costly. First, filter at the source rather than downloading everything and subsetting locally, because the QCEW single-county call and the \texttt{sub()} option on \texttt{educationdata} move only the rows you need across the network, which is lighter on the server and faster for you. Second, cache what you pull: save the raw download to \texttt{\$raw} and rerun your analysis against the local copy, so a debugging session that reruns a do-file forty times hits the disk forty times and the public endpoint once. Third, identify yourself where an API asks for a key or a contact, because an endpoint that can see who is calling can warn a heavy user before it has to block one. Each move protects the commons and, not by accident, makes the download replicable, since a cached raw file is the thing a colleague reruns a year later when the live endpoint has moved.

\section{Census data with getcensus}
The American Community Survey (ACS)\index{American Community Survey (ACS)} is the evaluator's standing source for community context: income, poverty, insurance, education, and dozens more, at geographies from the nation down to the census tract. The community package \texttt{getcensus}\index{getcensus@\texttt{getcensus}} turns an ACS query into one Stata command. It needs a Census API key, which is free and takes a minute: request one at \texttt{api.census.gov/data/key\_signup.html}, and the Census Bureau emails it to you. Store it once in your \texttt{profile.do} as \texttt{global censuskey "YOUR\_KEY"} (the full setup is in Appendix~\ref{app:setup}) so it loads at startup and never lands in a shared do-file. With that key in your profile, a county income panel is four lines:
{\footnotesize
\begin{verbatim}
* one-time, in profile.do:  global censuskey "YOUR_KEY"
getcensus B19013, years(2019/2023) ///
    geography(county) statefips(48) clear
\end{verbatim}
}
The command \texttt{getcensus catalog, search(poverty)} searches the data dictionary so you can find the table code without leaving Stata.\marginnote{Watch the margins of error: every ACS estimate ships with one, and the five-year file exists precisely because a single year's sample is too thin for small geographies. Below the county level, report the estimate and its interval together, never the point alone.} Before an evaluator leans on the ACS for a board exhibit, the Census Bureau's own handbook for data users is worth an afternoon: it walks through period estimates, margins of error, and the geographies the survey supports, in language written for practitioners rather than statisticians \parencite{census2020acs}. Reading it once is what lets an embedded analyst answer a program director's \enquote{is this number solid?} without hedging, which is the accessible principle applied to the evaluator's own understanding, not just the reader's. One honest limit: \texttt{getcensus} covers the ACS, not the decennial census or CPS microdata; for those endpoints, use the Python bridge of Section~\ref{sec:threeroutes}. This is the section that supplies the benchmark data behind the rural-reach example of Chapter~\ref{ch:embedded}, and it serves accessibility directly, because it puts real geography and real units into a funder's report on demand.

\section{Education panels with educationdata}
The Urban Institute's Education Data Portal harmonizes federal education data (the Common Core of Data on schools and districts, IPEDS on colleges, and the Civil Rights Data Collection) so that variable names stay consistent across years, which is exactly the property that makes a panel painless to build. The portal's builders set out to solve the exact problem an embedded evaluator hits at the start of every education project: the raw federal files are public but scattered across formats and inconsistent codings, so the difficulty of assembling them slows evidence-based decisions to a crawl \parencite{chingos2018portal}. The portal absorbs that assembly cost once, on behalf of everyone, so the analyst inherits a clean panel rather than a cleaning project. The official \texttt{educationdata}\index{educationdata@\texttt{educationdata}} package needs no key:
{\footnotesize
\begin{verbatim}
educationdata using "school ccd enrollment", ///
    sub(year=2015:2022 fips=48) clear
educationdata using "school ccd directory", meta
\end{verbatim}
}
The \texttt{sub()} option subsets on the server so you download only what you need, and \texttt{meta} returns the codebook without downloading data at all.\marginnote{The portal wraps an open REST API, so the same query runs from R (\texttt{educationdata} on CRAN) or a browser URL. If a colleague works in R, you are pulling the identical rows, which is one less reconciliation meeting.} Someone else has already done the cross-year name harmonization that Chapter~\ref{ch:panels} treats as hard-won; here it arrives as a gift, and the panel lands analysis-ready. That is extensibility bought cheaply: add a year to \texttt{sub()} and the panel grows.

\subsection{A worked example: did Texas enrollment dip with COVID?}\label{sec:ccd}
A school-finance analyst in 2023 wants to know how far the pandemic pushed enrollment down, and by how much, because the district's per-pupil funding follows the headcount. The Common Core of Data has the answer at the district level for every year, and \texttt{educationdata} pulls the Texas panel with one filtered request. We ask the server for district totals only (\texttt{grade=99} is the all-grades line) across 2015 through 2022, so only the rows we need cross the wire:
{\footnotesize
\begin{verbatim}
educationdata using "district ccd enrollment", ///
    sub(year=2015:2022 grade=99 fips=48) clear
assert grade == 99
drop if missing(enrollment) | enrollment < 0
collapse (sum) enrollment (count) n_dist=enrollment, ///
    by(year) fast
list year enrollment n_dist, clean noobs
\end{verbatim}
}
The \texttt{collapse} sums the district counts into one state total per year and, in the same pass, counts how many districts reported, which is the denominator you want before trusting any state number. The result is an eight-row series, and the numbers are real:
{\footnotesize
\begin{verbatim}
    year    enrollment    n_dist
    2015     5,301,916      1226
    2016     5,360,847      1231
    2017     5,401,097      1236
    2018     5,433,738      1240
    2019     5,494,830      1244
    2020     5,373,203      1247
    2021     5,428,611      1250
    2022     5,519,724      1252
\end{verbatim}
}
Read this in students, not proportions. Texas gained roughly 48{,}000 students a year from 2015 to 2019, then lost about 122{,}000 between the fall of 2019 and the fall of 2020, a drop of 2.2 percent in a single year that erased more than two years of prior growth. By 2022 the count had not just recovered but passed the old peak, reaching 5.52 million. The \texttt{n\_dist} column climbs steadily from 1{,}226 to 1{,}252 reporting districts, so the dip is a real loss of students, not an artifact of districts dropping out of the file. This makes sense: the pandemic's enrollment shock hit in the 2020 fall count and unwound over the next two years, the same V-shape that district finance offices across the state were bracing for.

\begin{figure}[htbp]
\centering
{\footnotesize
\begin{verbatim}
twoway connected enroll_m year, ///
    xline(2020, lpattern(dash)) ///
    title("Texas public school enrollment," ///
          " 2015-2022")
\end{verbatim}
}
\includegraphics[width=\linewidth]{ch03_ccd_enrollment.png}
\caption[Texas public school enrollment, 2015--2022]{Texas public school enrollment, 2015--2022, built by \texttt{ch03\_educationdata.do} from the Common Core of Data through the \texttt{educationdata} package with no API key. The dashed line marks the fall of 2020: enrollment falls about 122{,}000 students from its 2019 peak, then recovers past it by 2022. Add a year to \texttt{sub()} and the series extends.}
\label{fig:ccd}
\end{figure}

The whole example, download to figure, is one short do-file (\texttt{ch03\_educationdata.do}) that reruns from the same portal call, so next year's data point arrives by editing one number in \texttt{sub()}. The cross-year name harmonization that Chapter~\ref{ch:panels} treats as hard-won has already been done for you on the server, so the panel lands analysis-ready. That is extensibility bought cheaply, and it hands a finance office a defensible enrollment trend it can act on rather than a rumor about the COVID year.

\section{Economic series with import fred}
Economic context, unemployment and wages, is what nearly every workforce or place-based evaluation needs for its comparison story. Stata reads the Federal Reserve's FRED\index{import fred@\texttt{import fred}}\index{FRED (Federal Reserve Economic Data)} database natively, no package required:
{\footnotesize
\begin{verbatim}
set fredkey YOUR_KEY, permanently
import fred TXUR CAUR NYUR, daterange(2010-01-01 .) clear
\end{verbatim}
}
FRED returns wide data, one column per series, so a \texttt{reshape long} is the first teaching moment in panel construction.\marginnote{FRED series are also mixed-frequency: \texttt{TXUR} is monthly but a GDP series is quarterly, and stacking them without aligning the date grid leaves ragged gaps. Decide your target frequency before you reshape, not after.} FRED needs a free key (Appendix~\ref{app:setup}); the next source needs none at all.

\section{A worked example: county wages from the BLS, no key}\label{sec:qcew}
A workforce board asks a concrete question: what does a job pay here, compared with the metros we compete with? The Bureau of Labor Statistics answers it through the Quarterly Census of Employment and Wages (QCEW)\index{Quarterly Census of Employment and Wages (QCEW)}, and it posts the data as plain CSV files, one per county-year, that Stata reads directly from a URL with no key.\marginnote{The QCEW counts jobs, not people: a worker with two jobs is counted twice, and it covers only employment subject to unemployment insurance, so the self-employed and most agricultural labor are excluded. For a headcount of residents, reach for the Census ACS instead.} Each file is downloaded by a single \texttt{import delimited}. To compare five Texas metros we loop over their county FIPS codes and keep the one row we want, the total private-sector line (ownership code~5, industry code~10):
{\footnotesize
\begin{verbatim}
* Travis, Harris, Dallas, Bexar, Tarrant:
local counties 48453 48201 48113 48029 48439
local base "https://data.bls.gov/cew/data/api/2023/1/area"
tempfile master
local first 1
foreach fips of local counties {
    import delimited ///
        "`base'/`fips'.csv", ///
        clear varnames(1) stringcols(1 3)
    keep if own_code == 5 & industry_code == "10"
    keep area_fips avg_wkly_wage month3_emplvl
    if `first' local first 0
    else       append using `master'
    save `master', replace
}
\end{verbatim}
}
One detail earns a comment. We force columns~1 and~3 (the county FIPS code\index{FIPS codes} and the industry code) to import as strings with \texttt{stringcols(1 3)}, because they are identifiers, not quantities: read as numbers, \texttt{"48453"} would keep its value but \texttt{"10"} could collide with other codes, and leading zeros in smaller FIPS codes would vanish. Reading identifiers as text is a habit that prevents a whole class of silent merge failures later.\marginnote{The classic failure: a five-digit FIPS like \texttt{01001} (Autauga, Alabama) loses its leading zero as a number, becomes \texttt{1001}, and fails to merge against a crosswalk that stored it as a string. Every FIPS below Colorado is at risk.}

After labeling each county for a client-readable chart, the extract is five rows, and the numbers are real:
{\footnotesize
\begin{verbatim}
    county                 month3_emplvl   avg_wkly_wage
    Harris (Houston)            2129826            1,900
    Travis (Austin)              660281            1,862
    Dallas                      1644115            1,839
    Tarrant (Ft. Worth)          872936            1,391
    Bexar (San Antonio)          764432            1,270
\end{verbatim}
}
This makes sense: Houston, Austin, and Dallas anchor the state's corporate and technology payrolls, while Fort Worth and San Antonio sit several hundred dollars a week lower, a gap of roughly a third. A workforce board reading this sees immediately that a training program aiming at the regional wage will look different in Bexar than in Harris. One \texttt{graph hbar} turns the extract into the exhibit:
{\footnotesize
\begin{verbatim}
graph hbar (asis) avg_wkly_wage, ///
    over(county, sort(1) descending label(labsize(small))) ///
    bar(1, color(navy)) blabel(bar, format(%9.0fc)) ///
    title("Private-sector weekly wage, Texas metros, 2023 Q1")
graph export "$figures/ch03_qcew_wages.png", ///
    replace width(2400)
\end{verbatim}
}

\begin{figure*}[htbp]
\centering
\includegraphics[width=\linewidth]{ch03_qcew_wages.png}
\caption[County wages from the BLS QCEW]{Average private-sector weekly wage across five Texas metro counties, 2023 Q1, built end to end by \texttt{20\_ch03\_apis.do} from public BLS data with no API key. The command that made this figure is printed above it; change the county codes or the year and it rebuilds.}
\label{fig:qcew}
\end{figure*}

The whole example, download to figure, is one short do-file (\texttt{20\_ch03\_apis.do}) that anyone can rerun from the same URLs. Loop the same code over more quarters and years and the five-row snapshot becomes a county-quarter panel, which is where Chapter~\ref{ch:panels} picks it up. Zero keys and a plain \texttt{import delimited} is as replicable as data acquisition gets, and the payoff is a comparison the board did not have before.

\section{Health and community context: County Health Rankings and PLACES}
Small-area health estimates make community need visible to a funder who thinks in counties. County Health Rankings, produced annually since 2010 by the University of Wisconsin Population Health Institute and the Robert Wood Johnson Foundation, ranks every county in a state against its peers on health outcomes and the factors behind them, and its authors describe it as a \enquote{population health checkup} built to spur local action rather than to settle a research debate \parencite{remington2015rankings}. That framing suits an embedded evaluator exactly: the file is designed to be acted on, and its model of health factors maps cleanly onto the community-level outcomes a logic model already names. County Health Rankings publishes an annual analytic file, one row per county, at a stable per-year URL, so a plain \texttt{copy} then \texttt{import} pulls a national county dataset of premature death, child poverty, and clinical-care measures with no key and no login:
{\footnotesize
\begin{verbatim}
local site "https://www.countyhealthrankings.org"
local path "sites/default/files/media/document"
copy "`site'/`path'/analytic_data2024.csv" ///
     "$raw/chr2024.csv", replace
import delimited "$raw/chr2024.csv", clear varnames(1)
\end{verbatim}
}
The download is a real one, roughly 3{,}200 counties and several hundred columns, and its header is the kind of thing you meet constantly: the file's first data row is a second label row, and the column names are long phrases like \texttt{Premature Death raw value} that Stata munges into \texttt{prematuredeathrawvalue}. Cleaning that (dropping the label row, renaming the handful of columns you need) is the harmonization work of Chapter~\ref{ch:panels}, previewed here so the download does not look tidier than it is.

The CDC's PLACES project is the other standing source, modeling roughly forty health measures down to the tract through Socrata\index{Socrata} endpoints that emit CSV directly. PLACES exists to put local prevalence estimates where almost none existed before: small-area survey samples are too thin to report chronic-disease rates for a single tract, so the project models them from national survey data down to four geographic levels, county through ZIP-code area, on one consistent method \parencite{greenlund2022places}. For an evaluator whose service area is a handful of neighborhoods, that modeled estimate is often the only defensible health number available below the county, which is why it is worth the extra care its Socrata endpoint demands.\marginnote{Socrata speaks SoQL, a SQL-like query language, so you can push \texttt{\$where} and \texttt{\$select} clauses server-side and download only the columns and rows you need, the same filter-at-the-source instinct that Chapter~\ref{ch:noapi} scales to 100\,GB files.} Socrata is worth knowing, and it carries two traps worth a boxed warning, one of them specific to Stata (see the box below).

\begin{vizcallout}[The Socrata truncation trap]
Without a row limit, a Socrata endpoint returns only the first 1{,}000 rows, so a state's worth of tracts arrives truncated but looks complete. Constrain the request with field filters (as in the \texttt{webapi} examples of Section~\ref{sec:webapi}), and confirm the count with \texttt{count} after import. Socrata's own \texttt{\$limit} option would do this too, but its leading dollar sign collides with Stata's macro syntax; Section~\ref{sec:webapi} shows the clean way around it. Silent truncation is the kind of error that survives all the way to a published table.
\end{vizcallout}

These public data sources make an evaluation actionable: they let the analyst turn \enquote{this community has real need} from an assertion into a number a board can act on.

\section{Three routes from Stata to any API}\label{sec:threeroutes}
Every API you will meet reduces to one of three routes, and knowing which one a new endpoint needs is what lets you reuse a template instead of learning each agency from scratch. The first route is a bare \texttt{import delimited} from a URL, which works whenever the endpoint returns comma-separated text and needs no login, as PLACES and QCEW do above. The second route is for endpoints that hand back JSON or want an authorization token, and it is the one that used to mean writing Python; the \texttt{webapi}\index{webapi@\texttt{webapi}} command (Section~\ref{sec:webapi}) now collapses it into a single line. The third route runs the other direction, writing data out to a web service, and it is how the \texttt{googlesheets}\index{googlesheets@\texttt{googlesheets}} and \texttt{googlechart}\index{googlechart@\texttt{googlechart}} tools of Chapters~\ref{ch:sheets} and~\ref{ch:interactive} publish results back to the browser.

Table~\ref{tab:ch03-viz-routes} is the whole decision on one line each: match the reply an endpoint gives you to a row, and you have the command and the template.

\begin{table}[htbp]
\centering
\footnotesize
\begin{tabular}{@{}p{1.4cm}p{1.7cm}p{2.0cm}p{3.2cm}@{}}
\toprule
Route & Reply looks like & Command & Reach for it when \\
\midrule
Flat URL & CSV, no login & \texttt{import}\newline\texttt{delimited} & the endpoint hands back plain comma-separated text and needs no key (QCEW, PLACES) \\
\addlinespace
\texttt{webapi} & nested JSON, or a token wall & \texttt{webapi get} & the reply is nested JSON, or the endpoint wants a bearer token or an API key \\
\addlinespace
Write-out & you send data up & \texttt{googlesheets}, \texttt{googlechart} & you are publishing results back to the browser, not pulling them down \\
\bottomrule
\end{tabular}
\caption[Three routes from Stata to any API]{Which route a new endpoint needs, decided by the reply it returns. The first two pull data down; the third writes it out. Identify the row and you reuse a template instead of learning the agency from scratch.}
\label{tab:ch03-viz-routes}
\end{table}

Owning one template per route is what makes this skill extensible past the agencies named here. When a state agency stands up a new endpoint next year, you will not need a new chapter; you will recognize which route applies and reuse it. That is the difference between learning five APIs and learning how APIs work.

\subsection{JSON without tears, and the \texttt{webapi} command}\label{sec:webapi}
The moment an evaluator hits an API that answers in JSON rather than a tidy CSV, the twenty-minute job threatens to become a Python project. This section is the shortcut: what JSON is, in one paragraph, and the one command that turns it into a Stata dataset.

Start with the format, because the word scares people more than the thing does. \emph{JSON}\index{JSON (JavaScript Object Notation)} (JavaScript Object Notation) is just text that stores data as labeled boxes inside labeled boxes. A record for one person might read \texttt{\{"id":1, "name":"Leanne Graham", "address":\{"city":"Gwenborough"\}\}}: three fields, and the third, \texttt{address}, is itself a small box holding \texttt{city}. That nesting is the only thing that makes JSON harder to read than a spreadsheet, because a spreadsheet cannot put a box inside a cell. The fix is \emph{flattening}\index{flattening (JSON)}: turning \texttt{address.city} into a plain column named \texttt{addresscity}, so the nested box becomes an ordinary variable.

The \texttt{webapi} command does the request, the authentication, and the flattening in one line, and it leans only on the Python that ships inside Stata, so there is nothing to \texttt{pip install}. Ask a public endpoint for its data and \texttt{webapi} hands you a dataset:
{\footnotesize
\begin{verbatim}
webapi get using ///
    "https://jsonplaceholder.typicode.com/users", clear
list id name addresscity in 1/3, noobs
\end{verbatim}
}
The reply is ten records with nested addresses, and the flattening has already happened, so \texttt{address.city} arrives as the column \texttt{addresscity} with no work on your part:
{\footnotesize
\begin{verbatim}
  +------------------------------------------+
  | id               name        addresscity |
  |------------------------------------------|
  |  1      Leanne Graham        Gwenborough |
  |  2       Ervin Howell        Wisokyburgh |
  |  3   Clementine Bauch      McKenziehaven |
  +------------------------------------------+
\end{verbatim}
}
That is the whole idea: you name the URL, and the nested reply lands as flat variables. This makes sense once you see it, and it is why an evaluator can treat a JSON API as just another data source rather than a programming detour.

The same one line reaches real health data. The CDC serves national mortality through a JSON endpoint, and a field filter keeps the request to the rows you want, still with no key:
{\footnotesize
\begin{verbatim}
webapi get using ///
    "https://data.cdc.gov/resource/bi63-dtpu.json", ///
    params("year=2017|cause_name=All causes") clear
\end{verbatim}
}
That returns 52 rows (one per state, plus the District of Columbia and the national total) over HTTP~200, ready to \texttt{destring} and rank. The \texttt{params()} option takes pipe-separated \texttt{key=value} pairs and encodes them into the query string for you, so you never hand-build a URL.

When an endpoint does need a key, \texttt{webapi} carries it without ever putting it in your do-file. Store the token once as a global in \texttt{profile.do}, then pass it by reference:
{\footnotesize
\begin{verbatim}
* one-time, in profile.do:  global API_TOKEN "your_token"
webapi get using "https://api.example.org/v1/series", ///
    bearer("$API_TOKEN") params("since=2026-01-01") ///
    records("data") clear
\end{verbatim}
}
Three options cover the auth schemes you will meet: \texttt{bearer()} for a token, \texttt{basicauth("user:pw")} for a username and password, and \texttt{apikeyheader()} or \texttt{apikeyparam()} for a plain API key. The \texttt{records("data")} option points at the branch of the reply that holds the rows, for the common case where an API wraps its data in an envelope of metadata. Keeping the key in \texttt{profile.do} rather than the script is the same discipline as every other credential in this book: a key in a shared file is a key someone else will spend.

\begin{vizcallout}[The one JSON trap in Stata: the \$ sign]
Socrata endpoints (data.cdc.gov and many state portals) offer query options that begin with a dollar sign, such as \texttt{\$limit} and \texttt{\$where}. Stata reads \texttt{\$} as the start of a global macro, so \texttt{params("\$limit=5000")} silently collapses to \texttt{params("=5000")} and the request fails. The clean fix is to not use the dollar options at all: the plain field filters shown above (\texttt{field=value}) do the same filtering with no dollar sign, and if you only need to cap the row count you can \texttt{keep in 1/5000} in Stata after the import. It is the one non-obvious trap in the chapter, and it bites everyone once.
\end{vizcallout}

\begin{kaobox}[frametitle=Package Integration: \texttt{webapi}]
\textbf{Integrated command:} \texttt{webapi get}/\texttt{post}. Fetches a JSON (or CSV) web reply, optionally with a bearer token, basic auth, or an API key, and flattens the nested reply into a Stata dataset, using only Stata's built-in Python. It also polls an endpoint on a timer (\texttt{every()}, \texttt{times()}) to build a live panel, which is the survey-monitoring pattern of Chapter~\ref{ch:surveys}. Install it once from the authors' GitHub:
{\footnotesize
\begin{verbatim}
local gh "https://raw.githubusercontent.com/ericabooth"
net install webapi, ///
    from("`gh'/webapi-stata-public/main/") replace
\end{verbatim}
}
The companion \texttt{ch03\_webapi.do} runs every example here against live public endpoints with no key.
\end{kaobox}

This matters because it widens an evaluator's reach: once JSON is no harder than CSV, an embedded evaluator on a deadline can pull data from far more than the handful of agencies that publish flat files, reaching the far larger set that serve JSON, which is most modern APIs. That is the accessible and extensible principle pointed at the whole open-data landscape, not just the friendly corner of it.

\begin{kaobox}[frametitle=Package Integration: \texttt{combineall}]
\textbf{Integrated command:} \texttt{combineall}\index{combineall@\texttt{combineall}}. The \enquote{combine} half of the download loops above: point it at a folder and it appends (or merges, or converts) every file inside, and with a rename-map CSV (old name, new name, first year, last year) it applies vintage-aware renames, stamps each variable's source into a characteristic, and reports what failed to match, so a multi-year panel becomes one command instead of a fragile copy-and-edit. Chapter~\ref{ch:noapi} works it in full.
\end{kaobox}

\begin{appliedexample}[A Texas county panel from three APIs]
\textbf{Scenario.} A county workforce board wants a one-page context exhibit: income, child poverty, and unemployment for its counties, 2015--2023, next to the state.

\medskip
\textbf{The build.} Pull ACS income and poverty with \texttt{getcensus} (county, \texttt{statefips(48)}); pull county unemployment from FRED with \texttt{import fred}; assemble QCEW wages from the no-key CSV loop. Stack the three with \texttt{combineall}, tagging each variable's source so the exhibit's provenance is auditable. The whole exhibit rebuilds next year by rerunning the do-file, which is why it is replicable, and it reports in the counties and dollars the board plans in, which is why it is actionable.
\end{appliedexample}

\paragraph{Where this leaves you.} You can now pull public data from the Census, education, economic, and health APIs into Stata, handle a JSON reply in one line with \texttt{webapi}, and place any new endpoint into the three-route pattern. Every one of those downloads is replicable because it runs from a URL rather than a screenshot, and JSON is no longer a reason to reach for a programmer. But plenty of agencies publish data with no API at all; Chapter~\ref{ch:noapi} is how you get it anyway.

\section*{Exercises}\addcontentsline{toc}{section}{Exercises}
\begin{enumerate}
  \item \emph{Try it on your own data.} Take a dataset your program already tracks that carries a county or state FIPS identifier, import that identifier column as a string, and confirm no leading zeros were lost by listing the shortest codes and checking that every value has the width you expect.
  \item Rerun \texttt{20\_ch03\_apis.do} for two more Texas counties by adding their FIPS codes to the \texttt{counties} local, and report where each new county's private-sector weekly wage falls relative to the five metros already in the exhibit.
  \item Predict, before you run anything, how many rows the CDC mortality call in Section~\ref{sec:webapi} returns for \texttt{year=2017} and \texttt{cause\_name=All causes}; then run the \texttt{webapi get} line and compare your guess to the count the chapter reports.
  \item Break the enrollment example on purpose: change the \texttt{assert grade == 99} line to \texttt{assert grade == 9} before the \texttt{collapse}, rerun it, and confirm the assertion stops the do-file rather than letting a wrong-grade total reach your report.
  \item Extend the FRED example by adding one more state's unemployment series to the \texttt{import fred} call, then \texttt{reshape long} the wide result and describe how the reshape aligns the three series onto one date grid.
\end{enumerate}

